\documentclass[aspectratio=169]{beamer}

% Unicode engine (LuaLaTeX or XeLaTeX) -> NO inputenc/fontenc
\usepackage[main=bulgarian]{babel}
\usepackage{tikz}

% --- НОВ ШРИФТ: PT Sans (Модерна кирилица) ---
\usepackage{fontspec}
\setsansfont{PT Sans}   % beamer is sans by default
\setmainfont{PT Sans}

\renewcommand{\familydefault}{\sfdefault}

% --- Полезни пакети ---
\usepackage{hyperref}
\usepackage{microtype}
\usepackage{amsmath,amssymb}
\usepackage{graphicx}

% --- ТЕМА И ЦВЕТОВЕ (Обновени) ---
\usetheme{Madrid}
% Използваме whale като солидна база за нашите цветове
\usecolortheme{whale} 
\setbeamertemplate{navigation symbols}{}

% Дефиниране на новите цветове
\definecolor{mainblue}{RGB}{0, 43, 92}    % Тъмносиньо
\definecolor{accentblue}{RGB}{0, 102, 204} % Акцентно синьо
\definecolor{softblue}{RGB}{230, 235, 245} % Светлосиньото, което заменя лилавото

% Прилагане на цветовете върху темата
\setbeamercolor{structure}{fg=mainblue}
\setbeamercolor{title}{fg=mainblue, bg=softblue}
\setbeamercolor{subtitle}{fg=mainblue!80, bg=softblue}
\setbeamercolor{frametitle}{fg=mainblue, bg=softblue}

% Настройка на блоковете (за научните ръководители и темата)
\setbeamercolor{block title}{fg=white, bg=mainblue}
\setbeamercolor{block body}{bg=mainblue!5}

% --- Полета и типография ---
\setbeamerfont{author in head/foot}{size=\tiny}
\setbeamersize{text margin left=10pt, text margin right=10pt}
\usefonttheme{professionalfonts}

\setbeamerfont{title}{series=\bfseries,size=\LARGE}
\setbeamerfont{subtitle}{series=\mdseries,size=\large}
\setbeamerfont{frametitle}{series=\bfseries,size=\Large}
\setbeamerfont{block title}{series=\bfseries}
\setbeamerfont{block body}{size=\normalsize}

% --- Блокове (чист дизайн със заоблени ъгли) ---
\setbeamertemplate{blocks}[rounded][shadow=true]

% --- Списъци ---
\setlength{\parskip}{2pt}
\setlength{\itemsep}{4pt}
\setlength{\topsep}{3pt}

% --- Линкове ---
\hypersetup{
  colorlinks=true,
  linkcolor=black,
  urlcolor=accentblue
}


% --- Заглавни данни ---
\title[Защита на дисертационен труд]{Защита на дисертационен труд}
\subtitle{Задочен докторант Ангел Иванов Димитриев}
\author{Институт по биофизика и биомедицинско инженерство}
\institute{ИБФБМИ–БАН \;|\; ФМИ към СУ „Св. Климент Охридски“}
\date{03.07.2026 г.}

\setbeamertemplate{footline}
{
  \leavevmode%
  \hbox{%
  
  % --- Институт ---
  \begin{beamercolorbox}[wd=.57\textwidth,ht=2.25ex,dp=1ex,left]{author in head/foot}%
    \usebeamerfont{author in head/foot}%
    \hspace*{2.5ex}\insertauthor%
  \end{beamercolorbox}%
  
  % --- Заглавие ---
  \begin{beamercolorbox}[wd=.23\textwidth,ht=2.25ex,dp=1ex,center]{title in head/foot}%
    \usebeamerfont{title in head/foot}\inserttitle%
  \end{beamercolorbox}%
  
  % --- Дата + номер ---
  \begin{beamercolorbox}[wd=.20\textwidth,ht=2.25ex,dp=1ex,center]{date in head/foot}%
    \usebeamerfont{date in head/foot}%
    \insertdate\hspace{2.5ex}%
    \insertframenumber{} / \inserttotalframenumber%
  \end{beamercolorbox}%
  
  }%
  \vskip0pt%
}

\begin{document}

% ------------------------------------------------------------------
% --- Слайд 1 (Title)
\begin{frame}[plain]
  \titlepage
  \vspace{-1.0em}
  \begin{block}{Научни ръководители}
    \begin{itemize}
      \item акад. дмн дтн Красимир Атанасов (ИБФБМИ–БАН)
      \item доц. д-р Нора Ангелова (ФМИ към СУ „Св. Климент Охридски“)
    \end{itemize}
  \end{block}
\end{frame}

% ------------------------------------------------------------------
% --- Слайд 2 - Тема
\begin{frame}{Тема на дисертационния труд}
    \begin{center}
        \begin{minipage}{0.8\textwidth}
            \begin{block}{.} 
                \centering
                \vspace{0.2cm}
                \large
                Софтуерен продукт за реализация на\\ 
                обобщеномрежови модели и негови приложения
                \vspace{0.2cm}
            \end{block}
        \end{minipage}
    \end{center}
\end{frame}

\begin{frame}{Структура на дисертационния труд}
    \begin{block}{Основни части на дисертационния труд}
        \vspace{0.15cm}

        \begin{itemize}
            \setlength{\itemsep}{0.9em}

            \item Увод
            \item \textbf{Глава 1:} Обобщени мрежи - теоретични основи и софтуер за симулации
            \item \textbf{Глава 2:} Нововъведени алгоритми за обобщени мрежи
            \item \textbf{Глава 3:} \textit{OnlineGN}: симулатор за обобщеномрежови модели
            \item \textbf{Глава 4:} Приложения на обобщените мрежи и симулатора \textit{OnlineGN}
            \item Заключение
            \item Научни приноси, публикации и библиография
        \end{itemize}

        \vspace{0.15cm}
    \end{block}
\end{frame}

% ------------------------------------------------------------------
% --- Слайд 3 - Недостатъци
\begin{frame}{Анализ на съществуващите софтуерни решения}
  \textbf{Основни недостатъци на текущите симулатори:}
  \vspace{0.2cm}

  \begin{columns}[T, totalwidth=\textwidth]
    \column{0.48\textwidth}
      \begin{block}{Неефективен процес на моделиране}
        \small 
        Ръчна работа при моделиране                                     и визуализация (извън \TeX).
      \end{block}

      \begin{block}{Трудна преносимост}
        \small 
        Проблеми с пренасяне/обмен на мрежи като файлове.
      \end{block}

    \column{0.48\textwidth}
      \begin{block}{Технически бариери}
        \small 
        Необходимост от инсталация и допълнителни зависимости.
      \end{block}

      \begin{block}{Функционален дефицит}
        \small 
        Ограничени възможности и липса на ключови функционалности за реални приложения.
      \end{block}
  \end{columns}
\end{frame}

% ------------------------------------------------------------------
% --- Слайд 4
\begin{frame}{Цел на дисертационния труд}
  
  \begin{block}{Цел}
    \centering \vspace{0.2cm} \large
    Разработване на уеб базиран онлайн симулатор за обобщени мрежи --- \textbf{\textit{OnlineGN}}.
    \vspace{0.2cm}
  \end{block}

  \vspace{0.4cm}

  \begin{columns}[t]
    \column{0.5\textwidth}
    \textbf{\color{mainblue}Среда и функционалност:}
    \begin{itemize}
      \item Интерактивна среда за създаване и визуализация.
      \item Редактиране, споделяне и симулация на модели.
    \end{itemize}

    \column{0.5\textwidth}
    \textbf{\color{mainblue}Валидация на приложимостта:}
    \begin{itemize}
      \item Разработване на реални ОМ модели.
      \item Симулация и верификация на моделите в системата.
    \end{itemize}
  \end{columns}
\end{frame}

% ------------------------------------------------------------------
% --- Слайд 5 - Задачи (НОМЕРИРАНИ)
\begin{frame}{Цел и задачи на дисертационния труд}
  \begin{columns}[T]
    % Лява колона
    \column{0.48\textwidth}
      \begin{block}{Задача 1}
        \vbox to 2cm{
          \vspace{0.1cm}
          \small Разработване и имплементация на алгоритъм за автоматично изчертаване на обобщени мрежи от структуриран текстов формат.
          \vfil
        }
      \end{block}

      \begin{block}{Задача 3}
        \vbox to 2cm{
          \vspace{0.1cm}
          \small Формулиране на изисквания към уеб базирано интегрирано решение.
          \vfil
        }
      \end{block}

    % Дясна колона
    \column{0.48\textwidth}
          \begin{block}{Задача 2}
        \vbox to 2cm{
          \vspace{0.1cm}
          \small Разработване и реализиране на алгоритъм за преобразуване на \textit{SVG} изображение на обобщена мрежа в \TeX{} формат за лесно използване в научни публикации.
          \vfil
        }
      \end{block}

      \begin{block}{Задача 4}
        \vbox to 2cm{
          \vspace{0.1cm}
          \small Проектиране на архитектурата на системата \textit{OnlineGN} и определяне на основните модули.
          \vfil
        }
      \end{block}
    \end{columns}
\end{frame}

% ------------------------------------------------------------------
% --- Слайд 6 - Задачи (НОМЕРИРАНИ)
\begin{frame}{Цел и задачи на дисертационния труд}
  \vspace{-0.25cm}
  \begin{columns}[T]
    % Горна лява колона
    \column{0.48\textwidth}
      \begin{block}{Задача 5}
        \vbox to 1.4cm{
          \vspace{0.1cm}
          \small Реализация на уеб базиран интерфейс за визуализация, редакция и персонализиране на генерираното изображение.
          \vfil
        }
      \end{block}

      \begin{block}{Задача 7}
        \vbox to 1.4cm{
          \vspace{0.1cm}
          \small Реализация на механизъм за онлайн съхранение и споделяне на мрежи чрез уникален линк в база данни.
          \vfil
        }
      \end{block}

    % Горна дясна колона
    \column{0.48\textwidth}
      \begin{block}{Задача 6}
        \vbox to 1.4cm{
          \vspace{0.1cm}
          \small Реализация на удобен интерфейс за редакция на предикати в индексираната матрица и характеристични функции.
          \vfil
        }
      \end{block}

      \begin{block}{Задача 8}
        \vbox to 1.4cm{
          \vspace{0.1cm}
          \small Създаване и формализиране на нови ОМ модели за реални процеси, непредставяни досега чрез обобщени мрежи.
          \vfil
        }
      \end{block}
  \end{columns}

  \vspace{0.2cm}

  % Долен широк блок за Задача 9
  \begin{block}{Задача 9}
    \centering
    \vspace{0.1cm}
    \small Симулация и верификация на моделите в \textit{OnlineGN} и оценка на точността и използваемостта на системата.
    \vspace{0.1cm}
  \end{block}
\end{frame}

% ------------------------------------------------------------------
% --- Слайд 7 - (Задача 1)
\begin{frame}{Алгоритъм за автоматично изчертаване на обобщени мрежи \\(Задача 1)}
    \begin{tikzpicture}[remember picture, overlay]
        
        % 1. JSON код
        \node (img1) at ([xshift=-3.5cm, yshift=-0.5cm]current page.center) {
            \includegraphics[height=0.75\textheight, keepaspectratio]{json.png}
        };

        % 2. Мрежа 
        \node (img2) at ([xshift=3.5cm, yshift=-0.5cm]current page.center) {
            \includegraphics[width=0.45\textwidth, keepaspectratio]{drawing.png}
        };

        % 3. Стрелка
        \draw[->, thick] ([xshift=0.7cm]img1.east) -- ([xshift=-0.1cm]img2.west);
            
    \end{tikzpicture}
\end{frame}

% ------------------------------------------------------------------
% --- Слайд 7.2 - (Задача 1)
\begin{frame}{Алгоритъм за автоматично изчертаване на обобщени мрежи \\(решение на Задача 1)}
    \begin{center}
        \begin{minipage}{0.87\textwidth}
            \begin{block}{Основни стъпки в алгоритъма}
                \vbox to 4.3cm{
                    \vspace{0.2cm}
                    \begin{itemize}
                        \setlength{\itemsep}{10pt}
                        \item Представяне на обобщената мрежа (ОМ) като ориентиран граф.
                        \item Създаване на текстова спецификация на графа с необходимите параметри за визуализация чрез Graphviz.
                        \item Създаване на изображение на графа чрез Graphviz.
                        \item Последваща обработка на полученото изображение с цел приближаване до конвенцията за графично представяне на ОМ.
                    \end{itemize}
                    \vfil
                }
            \end{block}
        \end{minipage}
    \end{center}
\end{frame}

% ------------------------------------------------------------------
% --- Слайд 8 - (Задача 2)
\begin{frame}{Алгоритъм за преобразуване на \textit{SVG} изображение на обобщена мрежа в \TeX{} формат (Задача 2)}
\begin{center}
\begin{tikzpicture}[baseline]
    \node (img1) at (-1,0) {\includegraphics[width=0.40\textwidth]{drawing.png}};
    \node (img2) at (6,0) {\includegraphics[width=0.40\textwidth]{converter.png}};
    \draw[->, thick] (img1.east) -- (img2.west);
\end{tikzpicture}
\end{center}
\end{frame}

% --- Слайд 8.2 - (Задача 2)
\begin{frame}{Алгоритъм за преобразуване на \textit{SVG} изображение на обобщена мрежа в \TeX{} формат (решение на Задача 2)}
    \begin{center}
        \begin{minipage}{0.75\textwidth}
            \begin{block}{Основни стъпки в алгоритъма}
                \vbox to 4.1cm{
                    \vspace{0.2cm}
                    \begin{itemize}
                        \setlength{\itemsep}{10pt}
                        \item Прилагане на глобални преобразувания.
                        \item Нормализация към координатна система на \textit{picture}.
                        \item Уеднаквяване на визуалните параметри.
                        \item Корекция на краищата на ребрата.
                        \item Създаване на \TeX{} изображение.
                    \end{itemize}
                    \vfil
                }
            \end{block}
        \end{minipage}
    \end{center}
\end{frame}

% ------------------------------------------------------------------
% --- Слайд 9 - (Задача 3)
\begin{frame}{Изисквания към уеб базираното решение (Задача 3)}
    \begin{columns}[T]
        % Функционални изисквания
        \column{0.45\textwidth}
        \begin{block}{Функционални изисквания}
            \vbox to 2.7cm{ 
                \vspace{0.2cm}
                \begin{itemize}
                    \item Визуализация на ОМ
                    \item Симулация на ОМ процес
                    \item Графично редактиране
                    \item Запис във файл
                \end{itemize}
                \vfil
            }
        \end{block}

        % Нефункционални изисквания
        \column{0.45\textwidth}
        \begin{block}{Нефункционални изисквания}
            \vbox to 2.7cm{
                \vspace{0.2cm}
                \begin{itemize}
                    \item Висока производителност
                    \item Висока надеждност
                    \item Комуникационна сигурност
                    \item GDPR съответствие
                \end{itemize}
                \vfil
            }
        \end{block}

    \end{columns}
\end{frame}


% ------------------------------------------------------------------
% --- Слайд 10 - (Задача 4)
\begin{frame}{Проектиране на архитектурата на системата \textit{OnlineGN} и определяне на основните модули (Задача 4)}
    \centering

    % Първа секция
    \textbf{Трислойна архитектура} \\
    \vspace{0.2cm}
    \includegraphics[width=0.5\textwidth]{arch-layers-overview.png} 
    
    \vspace{0.6cm} 

    % Втора секция
    \textbf{Топологично разположение и комуникация} \\
    \vspace{0.2cm}
    \includegraphics[width=0.5\textwidth]{cloud-deployment.png}
\end{frame}

% ------------------------------------------------------------------
% --- Слайд 11 - (Задача 4)
\begin{frame}{Проектиране на архитектурата на системата \textit{OnlineGN} и определяне на основните модули (Задача 4)}
    \textbf{Модули на \textit{OnlineGN}}:

    \setbeamercolor{block body example}{bg=mainblue!15, fg=black}
    \setbeamercolor{block title example}{bg=mainblue!30, fg=mainblue}

    \vspace{0.2cm}
    \centering
    \begin{minipage}{0.70\textwidth} 

        \begin{block}{}
            \centering \small Модул за устойчивост на данните (\textit{GenNet Persistence})
        \end{block}

        \begin{exampleblock}{}
            \centering \small Модул за графична визуализация (\textit{SVG/Graphviz})
        \end{exampleblock}

        \begin{block}{}
            \centering \small Модул за симулация на движението на ядрата
        \end{block}

        \begin{exampleblock}{}
            \centering \small Модул за импортиране и експортиране на ОМ
        \end{exampleblock}

        \begin{block}{}
            \centering \small Модул за визуализация на динамични процеси и анимация
        \end{block}

    \end{minipage}
\end{frame}

% ------------------------------------------------------------------
% --- Слайд 12 - (Задача 5)
\begin{frame}{Реализация на уеб базиран интерфейс (Задача 5)}
\begin{figure}[h]
    \centering
    \setlength{\fboxsep}{0pt}
    \setlength{\fboxrule}{1pt}
    {\color{mainblue}\fbox{\includegraphics[width=0.9\textwidth]{ui-main-screen.png}}}
\end{figure}
\end{frame}

% ------------------------------------------------------------------
% --- Слайд 13 - (Задача 6)
\begin{frame}{Реализация на удобен интерфейс за редакция на предикати в индексираната матрица и характеристични функции на позициите (Задача 6)}
    \begin{figure}[h]
        \centering
        \setlength{\fboxsep}{0pt}
        \setlength{\fboxrule}{1pt}
        {\color{mainblue}\fbox{\includegraphics[width=0.32\textwidth]{ui-edit-place-dialog.png}}}
    \end{figure}
\end{frame}

% ------------------------------------------------------------------
% --- Слайд 14 - (Задача 7)
\begin{frame}{Реализация на механизъм за онлайн съхранение и споделяне на мрежи чрез уникален линк в база данни (Задача 7)}
    \begin{figure}
        \centering
        \setlength{\fboxsep}{0pt}
        \setlength{\fboxrule}{1pt}
        
        \setbeamercolor{fboxcolor}{fg=mainblue}
        {\color{mainblue}
        
        \begin{minipage}{0.48\textwidth}
            \centering
            \fbox{\includegraphics[width=\textwidth]{ui-save-error.png}}
            \vspace{0.1cm}
            \small Неуспешен опит за съхранение на мрежа в базата от данни (пример за грешка).
        \end{minipage}
        \hfill
        \begin{minipage}{0.48\textwidth}
            \centering
            \fbox{\includegraphics[width=\textwidth]{ui-save-success.png}}
            \vspace{0.1cm}
            \small Успешно запазване на мрежа; генерирана хипервръзка на запазената мрежа.
        \end{minipage}
        }
    \end{figure}
\end{frame}

% ------------------------------------------------------------------
% --- Слайд 15 - (Задача 8 и Задача 9)
\begin{frame}{ОМ модели и симулации (Задача 8 и Задача 9)}
Представените модели са разработени с цел да адресират следните основни задачи:

    \begin{columns}[T]
        % Лява колона
        \column{0.48\textwidth}
            \begin{block}{Задача 8}
                \vbox to 1.5cm{
                    \vspace{0.1cm}
                    \small Създаване и формализиране на нови ОМ модели за реални процеси, непредставяни досега чрез обобщени мрежи.
                    \vfil
                }
            \end{block}

        % Дясна колона
        \column{0.48\textwidth}
            \begin{block}{Задача 9}
                \vbox to 1.5cm{
                    \vspace{0.1cm}
                    \small Симулация и верификация на моделите в \textit{OnlineGN} и оценка на точността и използваемостта на системата.
                    \vfil
                }
            \end{block}
    \end{columns}
\end{frame}

% ------------------------------------------------------------------
% --- Слайд 16
\begin{frame}{ОМ модели и симулации: Симулация на крайни автомати}
    \vspace{-0.2cm}
    \textbf{Идея на модела:}
    \vspace{0.2cm}

    \begin{block}{Логическо съответствие (Трансформация)}
        \centering
        \vspace{0.1cm}
        \large
        Състояния в автомат $\boldsymbol{\rightarrow}$ Преходии в ОМ \quad | \quad Преходи в автомат $\boldsymbol{\rightarrow}$ Позиции в ОМ
        \vspace{0.1cm}
    \end{block}

    \vspace{0.3cm}

    \begin{itemize}
        \setlength\itemsep{0.7em} 
        \item \textbf{Симулация:} Представяне на краен автомат чрез ОМ.
        \item \textbf{Паралелизъм:} Моделиране на недетерминизъм (НКА) чрез едновременно движение на паралелни ядра.
        \item \textbf{Критерий за успех:} Приемане на думата при достигане на финална позиция в мрежата.
        \item \textbf{Универсалност:} Единен подход за работа с ДКА и НКА.
    \end{itemize}
\end{frame}

% ------------------------------------------------------------------
% --- Слайд 17 
\begin{frame}{ОМ модели и симулации: Симулация на крайни автомати чрез ОМ }
Строене на ОМ по автомат.
\begin{center}
\begin{tikzpicture}[baseline]
    \node (img1) at (-1,0) {\includegraphics[width=0.40\textwidth]{aut.png}};
    \node (img2) at (6,0) {\includegraphics[width=0.40\textwidth]{OM-avtomat.png}};
    \draw[->, thick] (img1.east) -- (img2.west);
\end{tikzpicture}
\end{center}
\end{frame}

% ------------------------------------------------------------------
% --- Слайд 18
\begin{frame}{ОМ модели и симулации: Симулация на крайни автомати чрез обобщени мрежи }
    \begin{columns}[c] 
        \column{0.4\textwidth}
        \centering
        Начална конфигурация на ОМ, симулирана с \textit{OnlineGN}
        
        \column{0.6\textwidth}
        \centering
        \setlength{\fboxsep}{0pt}
        \setlength{\fboxrule}{1pt}
        {\color{mainblue}\fbox{\includegraphics[width=0.97\textwidth]{NFA-OGN-1.PNG}}}
    \end{columns}
\end{frame}

% ------------------------------------------------------------------
% --- Слайд 19 
\begin{frame}{ОМ модели и симулации: Паралелизация на задачата за Ханойските кули чрез обобщени мрежи}
    
    % Концептуална рамка
    \begin{block}{Идея на модела и разширение}
        \begin{itemize}
            \item Разширение на класическата задача
            \item Едновременно преместване на до $k$ пръстена
            \item Пръстените могат да се вземат от различни пръти
            \item \textbf{Цел:} по-малко стъпки чрез паралелизъм
        \end{itemize}
    \end{block}

    \vspace{0.1cm}

    % Правила и механизъм
    \begin{block}{Механизъм на симулацията}
        \begin{itemize}
            \item Едновременно повдигане на избраните пръстени
            \item \textbf{Поставяне по приоритет:} най-голям $\rightarrow$ най-малък
            \item Повдигане + поставяне = 1 паралелна стъпка
            \item \textbf{Запазва се правилото:} по-голям пръстен не стои върху по-малък
        \end{itemize}
    \end{block}
\end{frame}

% ------------------------------------------------------------------
% --- Слайд 20
\begin{frame}{ОМ модели и симулации: Паралелизация на задачата за Ханойските кули чрез обобщени мрежи}
    \begin{columns}[c] 
        \column{0.35\textwidth}
        \centering
        \begin{itemize} 
            \item $\rho$-ядро -- съответства на пръстен в задачата.
            \item $\chi$-ядро -- съответства на „ръка“ за преместване на пръстен.
        \end{itemize}

        \vspace{-0.2cm}
        \[
            \text{Брой стъпки} = 2^{\bigl\lceil \textstyle \frac{n}{k} \bigr\rceil} - 1
        \]

        \vspace{0.5cm}
        ОМ за Ханойските кули с $n=3$ $\rho$-ядра и $k=2$ $\chi$-ядра
        \vspace{0.5cm}
        \column{0.65\textwidth}
        \centering
        
        \includegraphics[width=0.48\textwidth]{TH OM1.png}
        \hspace{0.01\textwidth}
        \includegraphics[width=0.48\textwidth]{TH OM2.png}
    \end{columns}
\end{frame}

% ------------------------------------------------------------------
% --- Слайд 21 
\begin{frame}{ОМ модели и симулации: Паралелизация на задачата за Ханойските кули чрез обобщени мрежи }
\begin{columns}
    \column{0.4\textwidth}
    \hfill 
    \begin{minipage}{0.8\textwidth} 
        \centering
        ОМ за Ханойските кули с 3 \(\rho\)-ядра и 2 \(\chi\)-ядра в \textit{OnlineGN}
    \end{minipage}

    \column{0.6\textwidth}
    \centering
    \setlength{\fboxsep}{0pt}
    \setlength{\fboxrule}{1pt}
    
    {\color{mainblue}\fbox{\includegraphics[width=0.5\textwidth]{s1_1.png}}}
\end{columns}
\end{frame}

% ------------------------------------------------------------------
% --- Слайд 22
\begin{frame}{ОМ модели и симулации: Модел с ОМ на производството на газови и полимерни продукти в нефтопреработвателна рафинерия }
    \vspace{-0.2cm}
    \textbf{Идея на модела:}
    \vspace{0.3cm}

    \begin{block}{Обхват на симулацията}
        \begin{itemize}
            \item \textbf{Продукти:} Газово гориво, втечнен петролен газ, пропилен и полипропилен
            \item Моделиране на пълния производствен цикъл на газови и полимерни продукти в рафинерията
        \end{itemize}
    \end{block}

    \vspace{0.4cm}

    \begin{block}{Архитектура на ОМ модела}
        \begin{itemize}
            \item \textbf{Структура:} Преходите съответстват на инсталации в рефинерията
            \item \textbf{Логика:} Ядрата имат характеристика - текущия продукт и съответното количество. 
        \end{itemize}
    \end{block}
\end{frame}

% ------------------------------------------------------------------
% --- Слайд 23
\begin{frame}{ОМ модели и симулации: Модел с ОМ на производството на газови и полимерни продукти в нефтопреработвателна рафинерия}
    \begin{columns}[c] 
        \column{0.3\textwidth}
        \centering
        {Технологична схема}

        \column{0.7\textwidth}
        \centering
        \setlength{\fboxsep}{0pt}
        \setlength{\fboxrule}{1pt}
        
        {\color{mainblue}\fbox{\includegraphics[width=0.95\textwidth]{Figure1.png}}}
    \end{columns}
\end{frame}


% ------------------------------------------------------------------
% --- Слайд 24
\begin{frame}{ОМ модели и симулации: Модел с ОМ на производството на газови и полимерни продукти в нефтопреработвателна рафинерия }
\begin{columns}
\column{0.2\textwidth}
\centering 
ОМ модел

\column{0.8\textwidth}
\centering
\setlength{\fboxsep}{0pt}
\setlength{\fboxrule}{1pt}

{\color{mainblue}\fbox{\includegraphics[width=0.45\textwidth]{FG OM1.png}}}
\hspace{0.02\textwidth}
{\color{mainblue}\fbox{\includegraphics[width=0.45\textwidth]{FG OM2.png}}}
\end{columns}
\end{frame}

% ------------------------------------------------------------------
% --- Слайд 25
\begin{frame}{ОМ модели и симулации: Модел с ОМ на производството на газови и полимерни продукти в нефтопреработвателна рафинерия }
\begin{columns}
\column{0.2\textwidth}
\centering 
ОМ модел в \textit{OnlineGN}

\column{0.8\textwidth}
\centering

\setlength{\fboxsep}{0pt}  
\setlength{\fboxrule}{1pt}

{\color{mainblue}\fbox{\includegraphics[width=0.45\textwidth]{sim_2.png}}}
\hspace{0.03\textwidth}
{\color{mainblue}\fbox{\includegraphics[width=0.45\textwidth]{sim_1.png}}}
\end{columns}
\end{frame}

% ------------------------------------------------------------------
% --- Слайд 26
\begin{frame}{ОМ модели и симулации: Модел с ОМ на производството на тежки нефтопродукти в нефтопреработвателна рафинерия}
    \vspace{-0.2cm}
    \textbf{Идея на модела:}
    \vspace{0.3cm}
    \begin{block}{Обхват на симулацията}
        \begin{itemize}
            \item \textbf{Продукти:}
            \begin{itemize}
                \item Горива с масово съдържание на сяра \( \leq 0.5\% \), \( \leq 1.0\% \) и \( \leq 2.5\% \).
                \item Пътен битум: 50/70 и 70/100.
            \end{itemize}
            \item Моделиране на пълния производствен цикъл на тежки нефтопродукти в рафинерията
        \end{itemize}
    \end{block}

    \vspace{0.4cm}

    \begin{block}{Архитектура на ОМ модела}
        \begin{itemize}
            \item \textbf{Структура:} Преходите съответстват на инсталации в рефинерията
            \item \textbf{Логика:} Ядрата имат характеристика - текущия продукт и съответното количество. 
        \end{itemize}
    \end{block}

\end{frame}

% ------------------------------------------------------------------
% --- Слайд 27
\begin{frame}{ОМ модели и симулации: ОМ на производството на тежки нефтопродукти в нефтопреработвателна рафинерия}
    \begin{columns}[c] 
        \column{0.45\textwidth}
        \centering
        {Технологична схема}

        \column{0.55\textwidth}
        \centering
        \setlength{\fboxsep}{0pt}
        \setlength{\fboxrule}{1pt}

        {\color{mainblue}\fbox{\includegraphics[width=0.9\textwidth]{Figure_1.jpg}}}
    \end{columns}
\end{frame}

% ------------------------------------------------------------------
% --- Слайд 28
\begin{frame}{ОМ модели и симулации: Модел с ОМ на производството на тежки нефтопродукти в нефтопреработвателна рафинерия }
\begin{columns}
\column{0.15\textwidth}
\centering 
ОМ модел

\column{0.85\textwidth}
\centering
\setlength{\fboxsep}{0pt}
\setlength{\fboxrule}{1pt}

{\color{mainblue}\fbox{\includegraphics[width=0.48\textwidth]{OM FG3.png}}}
\hspace{0.02\textwidth}
{\color{mainblue}\fbox{\includegraphics[width=0.46\textwidth]{OM FG4.png}}}
\end{columns}
\end{frame}

% ------------------------------------------------------------------
% --- Слайд 29
\begin{frame}{ОМ модели и симулации: Модел с ОМ на производството на тежки нефтопродукти в нефтопреработвателна рафинерия }
\begin{columns}[c]
    \column{0.4\textwidth} 
    \centering 
    ОМ модел в \textit{OnlineGN}

    \column{0.6\textwidth}
    \centering
    \setlength{\fboxsep}{0pt}
    \setlength{\fboxrule}{1pt}

    {\color{mainblue}\fbox{\includegraphics[width=0.9\textwidth]{model5.png}}}
\end{columns}
\end{frame}

% ------------------------------------------------------------------
% --- Слайд 30
\begin{frame}{Научни приноси}
\scriptsize
\begin{enumerate}
    \setlength{\itemsep}{0.8em}

    \item \textbf{Алгоритми за автоматизация и визуализация на ОМ}
    \begin{itemize}
        \item Алгоритъм за автоматизирано изчертаване на ОМ по абстрактно описание.
        \item Алгоритъм за преобразуване на SVG изображения на ОМ в \TeX{} формат.
    \end{itemize}

    \item \textbf{Паралелен вариант на задачата за Ханойските кули}
    \begin{itemize}
        \item Формализиран е вариант с едновременно преместване на до $k$ пръстена.
        \item Доказана е теорема за минималния брой стъпки:
        \[
            H(n,k) = 2^{\bigl\lceil \textstyle \frac{n}{k} \bigr\rceil} - 1
        \]
    \end{itemize}

    \item \textbf{ОМ моделиране и симулация на паралелната задача}
    \begin{itemize}
        \item Създаден е оригинален ОМ модел с $\rho$-ядра за обектите и $\chi$-ядра за ресурсите.
        \item Приложимостта е доказана чрез симулации в \textit{OnlineGN}.
    \end{itemize}
\end{enumerate}
\end{frame}

% ------------------------------------------------------------------
% --- Слайд 31
\begin{frame}{Научно--приложни приноси}
\scriptsize
\begin{enumerate}
    \setlength{\itemsep}{0.8em}

    \item \textbf{Софтуерна система \textit{OnlineGN}}
    \begin{itemize}
        \item Реализирана е уеб базирана платформа за моделиране, симулация и визуализация на ОМ.
        \item Разработен е интерактивен интерфейс за редактиране и проследяване на изпълнението.
        \item Внедрен е експорт от SVG към \TeX{} за подготовка на публикации.
    \end{itemize}

    \item \textbf{Представяне на крайни автомати чрез ОМ}
    \begin{itemize}
        \item Разработена е методика за описание на детерминирани и недетерминирани крайни автомати чрез ОМ.
    \end{itemize}

    \item \textbf{Моделиране на процеси в нефтопреработвателна рафинерия}
    \begin{itemize}
        \item Разработени са ОМ модели на реални технологични процеси за газови, полимерни и тежки нефтени продукти.
    \end{itemize}

    \item \textbf{Верификация чрез \textit{OnlineGN}}
    \begin{itemize}
        \item Всички новосъздадени модели са верифицирани чрез симулации в \textit{OnlineGN}.
    \end{itemize}
\end{enumerate}
\end{frame}


% ------------------------------------------------------------------
% --- Слайд 30
\begin{frame}{Публикации}
\scriptsize
\begin{itemize}
    \setlength{\itemsep}{1.5em}

    \item Dimitriev, A.I.; Terziev, G.I. \textbf{Automating Creativity: Enhancing Generalized Nets with Algorithmic Drawing.} In: Atanassov, K.T., et al. (eds.) \textit{Uncertainty and Imprecision in Decision Making and Decision Support -- New Advances, Challenges, and Perspectives}. IWIFSGN 2024. \textit{Lecture Notes in Networks and Systems}, vol. 1857. Springer, Cham, 2026, pp. 62--75. doi:\,10.1007/978-3-032-22232-9\_7.

    \item Dimitriev, A. \textbf{Representing Finite-State Automata with Generalized Nets: Simulation via OnlineGN.} \textit{Proceedings of the Jangjeon Mathematical Society}, vol. 29, no. 1, 2026, pp. 155--167. doi:\,10.17777/pjms.2026.29.1.011.

    \item Dimitriev, A.; Atanassov, K.; Angelova, N. \textbf{Parallel Towers of Hanoi via Generalized Nets: Simulated with OnlineGN.} \textit{Software}, vol. 4, no. 4, article 23, 2025. doi:\,10.3390/software4040023.

    \item Stratiev, D.D.; Dimitriev, A.; Stratiev, D.; Atanassov, K. \textbf{Modeling the Production Process of Fuel Gas, LPG, Propylene, and Polypropylene in a Petroleum Refinery Using Generalized Nets.} \textit{Mathematics}, vol. 11, no. 17, article 3800, 2023. doi:\,10.3390/math11173800.

    \item Stratiev, D.; Dimitriev, A.; Stratiev, D.; Atanassov, K. \textbf{Generalized Net Model of Heavy Oil Products' Manufacturing in Petroleum Refinery.} \textit{Mathematics}, vol. 11, no. 23, article 4753, 2023. doi:\,10.3390/math11234753.
\end{itemize}
\end{frame}

% ------------------------------------------------------------------
% --- Слайд 31 - ДОПЪЛНЕН
\begin{frame}{Участия в научни конференции}
\scriptsize

\begin{block}{Участия, свързани с дисертационния труд}
\begin{itemize}
    \setlength{\itemsep}{0.7em}

    \item \textbf{22nd International Workshop on Intuitionistic Fuzzy Sets and Generalized Nets}
    (IWIFSGN'2024), Варшава, Полша -- 18.10.2024.

    Доклад: \textit{Automating Creativity: Enhancing Generalized Nets with Algorithmic Drawing}.

    \item \textbf{ИБФБМИ}, София, България -- 04.2024. Доклад: \textit{Представяне на OnlineGn}.

    \item \textbf{ИБФБМИ}, София, България -- 03.2025. Доклад: \textit{Представяне на OnlineGn}.
\end{itemize}
\end{block}

\vspace{0.3em}

\begin{block}{Друго научно участие}
\begin{itemize}
    \setlength{\itemsep}{0.5em}

    \item \textbf{International Symposium on Bioinformatics and Biomedicine}
    (BioInfoMed'2022), Бургас, България -- 05-07.10.2022.

    Доклади: \textit{ABO System Blood Groups Distribution in Bulgaria} и
    \textit{InterCriteria Analysis of the Geographic Distribution of the ABO System Blood Groups}.
\end{itemize}
\end{block}

\end{frame}

\end{document}